% MG05b. After inversion in the boundary fibre and before
% subsec:alpha-completa; requires tab:alpha-ventana and eq:alpha-telescopia.
\subsubsection{Controls of orientation, boundary, and register position}
\label{mg05:controles}

Integer closure distinguishes three operational dependencies:
the direction of carry propagation, the boundary condition, and the
position of the dodecaphasic register. Their comparison uses the same
twelve-block window as Table~\ref{tab:alpha-ventana}. The regional
words \(P,E,\Phi\) and the already generated register \(K\) are
retained except for the modification identified in each case. The
common reference is \(\mathcal C_{12}\) with \(c_{13}=0\), whose
output is \eqref{eq:alpha-ventana-cero}.

Prospective production of the signed register and its inversion
\(K^{(r)}=H_4U^{(r)}/4\) precede these controls. The normalization
of Lemma~\ref{lem:k-hadamard} is retained, in particular
\(H_4^2=4I_4\) and \(\det H_4=-16\). Inversion determines the
register once \(U\) has been produced; its operational selection
remains in the construction preceding closure. The following
variations of \(K\) examine the response of an already defined
reader.

\begin{proposition}[Separation of the dependencies of closure]
\label{mg05:dependencias}
For the data in Table~\ref{tab:alpha-ventana}, the following four
modifications have distinct effects:
\begin{enumerate}
\item Left-to-right quotient propagation, initialized with zero
quotient, changes the third block from \(352\) to \(353\).
\item The condition \(c_{13}=1\), with the original orientation,
changes the final block from \(663\) to \(664\) and preserves the
preceding eleven blocks.
\item The rotation \(K'_j=K_{j+1}\), with \(K_{13}=K_1\), changes
the output word while retaining \(c_{13}=0\).
\item The perturbation \(K'_1=K_1+1=235\), with the other eleven
coordinates fixed, changes the first block from \(007\) to \(006\)
and preserves the remaining blocks.
\end{enumerate}
\end{proposition}
\begin{proof}
For the first case, define the initial quotient \(b_0=0\) and
apply, in the order \(j=1,\ldots,12\),
\[
 \widetilde s_j=P_j+E_j-\Phi_j-K_j+b_{j-1},\qquad
 \widetilde A_j=[\widetilde s_j]_{1000},\qquad
 b_j=\frac{\widetilde s_j-\widetilde A_j}{1000}.
\]
This orientation produces, in two consecutive halves,
\begin{align*}
 \widetilde A_{1:6}&=(007,297,353,570,284,800),\\
 \widetilde A_{7:12}&=(995,285,106,471,379,663).
\end{align*}
The difference already appears at \(j=3\), where the quotient
required by the original closure comes from the fourth position.

For the second case, the final sum changes from \(663\) to \(664\).
Both have zero Euclidean quotient; hence \(c_{12}=0\) in both
cases and the recurrence agrees from \(j=11\) to \(j=1\).
The resulting word is
\begin{align*}
 A^{(1)}_{1:6}&=(007,297,352,569,283,800),\\
 A^{(1)}_{7:12}&=(997,285,105,472,380,664).
\end{align*}
Identity \eqref{eq:alpha-telescopia} records exactly the unit
variation of the right endpoint.

The rotation in the third case, evaluated with the same zero boundary
and the original recurrence, gives
\begin{align*}
 A^{\mathrm{rot}}_{1:6}&=(698,699,763,639,119,004),\\
 A^{\mathrm{rot}}_{7:12}&=(559,429,226,119,926,030).
\end{align*}
In this case the left carry is \(c_1=-1\); the right endpoint
remains zero. Retaining both endpoints in
\eqref{eq:alpha-telescopia} preserves the integer equality as the
twelve coefficients change position.

Finally, perturbing \(K_1\) leaves the entire recurrence from
\(j=12\) to \(j=2\) unchanged. At \(j=1\), the sum decreases
from seven to six, with zero quotient in both cases. This gives
\begin{align*}
 A^{\mathrm{pert}}_{1:6}&=(006,297,352,569,283,800),\\
 A^{\mathrm{pert}}_{7:12}&=(997,285,105,472,380,663).
\end{align*}
\end{proof}

The four controls separate transport direction, boundary fibre,
coefficient position, and local response. Their interpretation
retains the finite domain \(N=12\). Tail transport and normalization
of the complete section are developed next, with the compatibilities
required by passage to arbitrary depth.
